% Four-audit scope transfer: known algebraization input and proved complete-cover budgets.
\subsection{One henselian algebraization for all complete punctured covers}

\begin{lemma}[Known: fixed-base algebraization of the covering category]
\label{D:fixed-henselian-algebraization}
Let $(R,\mathfrak m,k)$ be a complete normal local surface ring
with algebraically closed coefficient and residue field $k$.
There is one normal strictly henselian local $k$-algebra $R^h$,
the henselization of a local ring essentially of finite type over
$k$, together with a fixed $k$-isomorphism
\[
                         \widehat {R^h}\simeq R.        \tag{D.AC.1}
\]
Completion induces an equivalence between the categories of finite
$R^h$-algebras etale away from the closed point and finite
$R$-algebras etale away from the closed point. It restricts to an
equivalence on normal local domains and preserves generic degree.
In particular the henselian base in this assertion is fixed before
any finite punctured cover is chosen.
\end{lemma}

\begin{proof}
The cited algebraization inputs are
\citet[Theorems~3.8--3.9 and~3.11, with Lemmas~3.12 and~3.17,
pp.~32--38]{Artin1969ApproximationD}. We explain the hypotheses
and the normal-domain restriction, rather than reprove those
known theorems.

Choose a power-series presentation $R=k[[X_1,\ldots,X_e]]/I$.
Over the perfect field $k$, the continuous Jacobian ideal has
support in the isolated singular locus. This is the same
formal Jacobian input used in the fixed-different construction
of Proposition~\ref{D:cover-uniform-canonical-defect}; its
continuous-differential interpretation is
\citet[Example~3.3, p.~7]{C_IyengarTakahashiStable}.
Thus the formal smoothness hypothesis outside the closed point
in Artin's Theorem~3.8 holds. That theorem gives an algebraic
local $k$-ring whose completion is $R$. Its residue field is
$k$, and its henselization $R^h$ has the same completion.
Faithfully flat descent of normality makes $R^h$ normal
\citep[Lemma~10.164.3, Tag~033G]{StacksCompleteBudgetD}.
Since $k$ is algebraically closed, $R^h$ is strictly henselian.
If $R$ is regular, the usual polynomial local ring and its
henselization provide the same conclusion directly.

Apply Artin's Theorem~3.11 to this fixed henselian ring and
its maximal ideal. It gives the asserted equivalence for
finite algebras, including all algebra homomorphisms. Let a
normal local cover domain $A$ over $R$ descend to a finite
algebra $A^h$ over $R^h$. The map
\[
              A^h\longrightarrow A^h\otimes_{R^h}R\simeq A
                                                               \tag{D.AC.2}
\]
is faithfully flat. Consequently $A^h$ is normal and is a
domain. A finite algebra over a henselian local ring is a
product of local rings; being a domain, $A^h$ is one local
ring. Its residue field is the finite extension $k$ of $k$.
The $\mathfrak m_{R^h}$-adic and maximal-ideal-adic topologies
on $A^h$ are cofinal, so (D.AC.2) identifies $A$ with its
local completion. Conversely, an excellent normal local
$A^h$ has normal local completion, hence a domain, and
completion preserves etaleness on the puncture. Generic
degree is preserved by flat base change: on the puncture
the finite algebra is locally free of its constant generic
rank, which is unchanged in (D.AC.2). This proves the
normal-domain and degree assertions.
\end{proof}

\begin{proposition}[Proved: degree-linear geometric budgets for all complete covers]
\label{D:all-complete-cover-geometric-budgets}
Let $R$ be as in Lemma~\ref{D:fixed-henselian-algebraization},
in characteristic $p>0$. For every finite normal local
domain $A/R$ of degree $t$ that is etale on the puncture,
the resolution-size, canonical-energy and integral-cycle
bounds of
Proposition~\ref{D:punctured-cover-linear-resolution},
Corollary~\ref{D:punctured-cover-linear-energy}, and
Corollary~\ref{D:punctured-cover-antinef-cycle}
hold with constants depending only on $R$ and fixed base
choices. No algebraizability hypothesis on the individual
complete cover needs to be added.

More explicitly, there is a regular projective resolution
with the degree-linear size bounds of the first cited
proposition. On the minimal regular resolution, with the
notation of the energy corollary,
\[
 v+\sum_i(-E_i^2)+\sum_i p_a(E_i)+p_g(A)
                +\kappa^{\mathsf t}N^{-1}\kappa=O(t).
                                                               \tag{D.AC.3}
\]
The empty exceptional set is interpreted as contributing
zero to the exceptional sums. There is also a full-support
anti-nef generation cycle $Z_0$ on that minimal resolution
such that
\[
 -Z_0\cdot B\ge2p_a(B)\quad(0<B\text{ exceptional}),
 \qquad -Z_0^2+|K\cdot Z_0|=O(t),                    \tag{D.AC.4}
\]
with the already specified fundamental-cycle treatment
of wholly zero-canonical configurations. When the minimal
exceptional set is empty take $Z_0=0$.

Consequently the additional energy hypothesis in
Proposition~\ref{D:hereditary-Cartier-generation-budget}
is satisfied for every such complete cover. On its
Cartier model one obtains the actual cycles
$W=\rho^*Z_0+C\ge C$ and all the energy, generation and
fixed-determinant finite-cycle chart bounds stated there,
uniformly in the degree $t$.
\end{proposition}

\begin{proof}
Fix the single $R^h$ of Lemma~\ref{D:fixed-henselian-algebraization}.
Every $A$ descends to one normal local $A^h/R^h$ of the
same degree. The strict-henselian clause of
Proposition~\ref{D:punctured-cover-linear-resolution}
already has constants determined by one fixed base
resolution and exceptional cycle. In particular its
constants do not change with the etale neighbourhood to
which a particular cover descends. Apply it to all the
$A^h$ over this fixed $R^h$.

We check precisely what completion preserves. Let
$Y^h\to\operatorname{Spec}A^h$ be any regular projective
resolution that is an isomorphism on the puncture and put
\[
                   Y=Y^h\times_{A^h}\operatorname{Spec}A.
                                                               \tag{D.AC.5}
\]
The complete local rings at corresponding points of the
closed fibres agree. In particular, a proper finite-type
scheme over a Noetherian local ring is regular if and
only if its base change to the completion is regular:
\citet[Lemmas~54.11.1--54.11.2, Tag~0BG4]
{StacksCompleteBudgetD}. Hence (D.AC.5) is a regular
projective resolution. It remains an isomorphism off the
closed fibre. The closed fibres, with their scheme
structures over $k$, are literally the same.

Each exceptional integral curve and each effective
exceptional cycle therefore corresponds to the same
proper $k$-scheme on both sides. To see the assertion
also for a thickened cycle, its structure sheaf is killed
by some power of $\mathfrak m_{A^h}$; the quotients of
$A^h$ and $A$ by that power agree. The normal line
bundles of the corresponding Cartier curves agree as
well. Thus all intersection numbers, arithmetic genera,
and cohomology of exceptional cycles are unchanged.
Canonical intersections are unchanged by adjunction,
$K\cdot E_i=2p_a(E_i)-2-E_i^2$; no assertion about
ordinary absolute differential forms after completion
is required.

Flat coherent base change gives
\[
 H^1(Y^h,\mathcal O_{Y^h})\otimes_{A^h}A
                    \simeq H^1(Y,\mathcal O_Y)
                                                               \tag{D.AC.6}
\]
by \citet[Lemma~30.5.2, Tag~02KH]{StacksCompleteBudgetD}.
The module on the left before tensoring has finite
length, so its $k$-dimension is unchanged. This proves
preservation of $p_g$. The same closed-fibre identification
shows that minimality of a regular resolution is
preserved: the exceptional smooth rational $(-1)$-curves
are exactly the same curves before and after completion.
One may equivalently complete a minimal resolution of
$A^h$. Therefore every entry of the intersection matrix,
the canonical vector and its energy is unchanged.

Apply the established strict-henselian size and energy
theorems and complete their resolutions. This proves
(D.AC.3) and the stated size assertion. The integral
anti-nef cycle construction and uniform-generation
criterion depend only on the same intersections,
arithmetic genera, and complete-cycle inequalities.
Their fixed sufficiently large choice of the core
positivity parameter yields (D.AC.4), with wholly
zero-canonical components treated by the already
proved rational-double-point fundamental cycle.
All constants remain those for the one fixed base
$R^h$. Finally apply
Proposition~\ref{D:hereditary-Cartier-generation-budget}.
This proves its conclusions for all the present
complete covers without altering its proof or silently
enlarging any earlier theorem's stated hypotheses.
\end{proof}

\begin{remark}[Scope of this transfer]
Artin's algebraization and finite-algebra comparison are
known inputs. The conclusion removes the former scope
qualification on the geometric budgets by using one
fixed henselian base for all covers. It does not produce
an arbitrary-family module or twisted-bundle atlas, a
uniform unit-Isom presentation, or a hereditary
$q^{Ctr^2}$ representation bound.
\end{remark}
